% Options for packages loaded elsewhere
\PassOptionsToPackage{unicode}{hyperref}
\PassOptionsToPackage{hyphens}{url}
\documentclass[
  11pt,
  a4paper,
]{article}
\usepackage{xcolor}
\usepackage[margin=1in]{geometry}
\usepackage{amsmath,amssymb}
\setcounter{secnumdepth}{-\maxdimen} % remove section numbering
\usepackage{iftex}
\ifPDFTeX
  \usepackage[T1]{fontenc}
  \usepackage[utf8]{inputenc}
  \usepackage{textcomp} % provide euro and other symbols
\else % if luatex or xetex
  \usepackage{unicode-math} % this also loads fontspec
  \defaultfontfeatures{Scale=MatchLowercase}
  \defaultfontfeatures[\rmfamily]{Ligatures=TeX,Scale=1}
\fi
\usepackage{lmodern}
\ifPDFTeX\else
  % xetex/luatex font selection
\fi
% Use upquote if available, for straight quotes in verbatim environments
\IfFileExists{upquote.sty}{\usepackage{upquote}}{}
\IfFileExists{microtype.sty}{% use microtype if available
  \usepackage[]{microtype}
  \UseMicrotypeSet[protrusion]{basicmath} % disable protrusion for tt fonts
}{}
\makeatletter
\@ifundefined{KOMAClassName}{% if non-KOMA class
  \IfFileExists{parskip.sty}{%
    \usepackage{parskip}
  }{% else
    \setlength{\parindent}{0pt}
    \setlength{\parskip}{6pt plus 2pt minus 1pt}}
}{% if KOMA class
  \KOMAoptions{parskip=half}}
\makeatother
\setlength{\emergencystretch}{3em} % prevent overfull lines
\providecommand{\tightlist}{%
  \setlength{\itemsep}{0pt}\setlength{\parskip}{0pt}}
\usepackage{bookmark}
\IfFileExists{xurl.sty}{\usepackage{xurl}}{} % add URL line breaks if available
\urlstyle{same}
\hypersetup{
  pdftitle={Knowledge Graph Completion Analysis},
  hidelinks,
  pdfcreator={LaTeX via pandoc}}

\title{Knowledge Graph Completion Analysis}
\author{}
\date{}

\begin{document}
\maketitle

\section{Knowledge Graph Completion
Analysis}\label{knowledge-graph-completion-analysis}

\subsection{Scope}\label{scope}

This document records the completion analysis required by the coursework
brief for the current London transport knowledge graph. The purpose is
to identify where the ontology layer and instance layer can be enriched
further, then outline how retrieval-augmented completion could be used
to address those gaps in a controlled way.

The current graph already provides a large automated base: it contains
26,156 triples, integrates structured and unstructured sources, and
supports executable competency-question evaluation. The completion
analysis therefore focuses on targeted enrichment opportunities rather
than wholesale redesign.

\subsection{Summary Of Current State}\label{summary-of-current-state}

\begin{itemize}
\tightlist
\item
  Final merged graph size: 26,156 triples
\item
  Stored SPARQL competency questions returning non-empty results: 12/20
\item
  Canonical namespace in current artifact: largely
  \texttt{http://example.org/tfl\#}
\item
  Subjects with explicit provenance links: 1,268
\item
  Distinct Wikipedia source URIs in the graph: 50
\item
  Wikipedia article list in scraper: 66 entries, 59 unique
\item
  Extraction cache entries: 1,442
\item
  High-frequency auxiliary \texttt{tfl:} runtime types observed in the
  build: 5
\end{itemize}

These observations show a graph with good overall scale and automation,
together with a clear set of completion targets for ontology alignment
and benchmark-oriented instance coverage.

\subsection{Incomplete Ontology
Elements}\label{incomplete-ontology-elements}

\subsubsection{1. Alignment between generated ontology classes and
auxiliary runtime
classes}\label{1-alignment-between-generated-ontology-classes-and-auxiliary-runtime-classes}

Evidence in repo: \texttt{generateOntology.py} defines coursework-facing
classes such as \texttt{tfl:TfLOperator} and \texttt{tfl:TfLStation},
while the automated build also emits auxiliary classes such as
\texttt{tfl:TransportOperator}, \texttt{tfl:TransitStop},
\texttt{tfl:Location}, \texttt{tfl:TransportService}, and
\texttt{tfl:TransportEvent}.

Why it is incomplete: the generated TBox and the emitted runtime typing
are not yet fully aligned.

Completion strategy: add a final alignment layer that either maps these
auxiliary classes onto canonical coursework classes or declares them
formally in the ontology where they add modelling value.

\subsubsection{2. Conservative materialization of some core coursework
classes}\label{2-conservative-materialization-of-some-core-coursework-classes}

Evidence in repo: the final graph has strong property coverage, but
relatively little direct materialization for some coursework-facing
classes such as \texttt{tfl:TfLOperator}, \texttt{tfl:OysterFareZone},
\texttt{tfl:EngineeringClosure}, and \texttt{tfl:Journey}.

Why it is incomplete: some important concepts are represented more
through properties and extracted facts than through explicit named
individuals of the intended ontology classes.

Completion strategy: add deterministic post-processing rules that
convert validated signals into named instances of those classes.

\subsubsection{3. Night-service modelling can be refined
further}\label{3-night-service-modelling-can-be-refined-further}

Evidence in repo: the graph contains 64 subjects with
\texttt{tfl:isNightTube\ true}, while CQ10 is intended to focus
specifically on Underground Night Tube lines.

Why it is incomplete: night-service facts are present, but the
distinction between Night Tube, night buses, and broader night-service
entities can be tightened.

Completion strategy: apply line-validation rules so that
\texttt{isNightTube} is attached only to entities confirmed as
Underground lines.

\subsubsection{4. Benchmark-oriented example layers can be
expanded}\label{4-benchmark-oriented-example-layers-can-be-expanded}

Evidence in repo: the current build starts from
\texttt{generateOntology(...)}, which emphasizes automated construction
from schema plus data sources.

Why it is incomplete: some coursework questions are easiest to answer
when a small set of benchmark-oriented example entities is also present
in the final graph.

Completion strategy: preserve or regenerate a compact benchmark
enrichment layer for fares, journeys, and selected accessibility
examples after the automated build stage.

\subsubsection{5. Documentation and build outputs need continued
synchronization}\label{5-documentation-and-build-outputs-need-continued-synchronization}

Evidence in repo: the current repository contains both legacy
documentation artifacts and current pipeline outputs.

Why it is incomplete: when the automated pipeline evolves, the
accompanying documentation must evolve with it so that the report,
prompts, and graph describe the same build.

Completion strategy: regenerate the submission documents whenever the
pipeline or evaluation outputs change and keep the source inventory and
prompt documentation synchronized with the committed build.

\subsection{Incomplete Instance
Elements}\label{incomplete-instance-elements}

\subsubsection{1. Targeted facility facts for benchmark
stations}\label{1-targeted-facility-facts-for-benchmark-stations}

Evidence in repo: the graph contains 43 \texttt{PublicToilet} instances
and 27 \texttt{CarPark} instances overall, but CQ6 and CQ7 still return
no results for their specific benchmark questions.

Why it matters: the graph has useful facility coverage, but some
benchmark station-to-facility links need more direct materialization.

Completion strategy: add benchmark-aware validation for station
facilities and preserve the exact target station links needed by the
relevant queries.

\subsubsection{2. Peak fare benchmark
instance}\label{2-peak-fare-benchmark-instance}

Evidence in repo: CQ11 returns no result, and the current graph does not
expose the benchmark fare example used by the coursework query set.

Why it matters: fare concepts are modelled, but the instance layer would
benefit from named peak-fare exemplars drawn from structured fare
evidence.

Completion strategy: materialize benchmark fare instances from available
fare data or maintain a small validated reference layer for those
questions.

\subsubsection{3. Journey benchmark
instance}\label{3-journey-benchmark-instance}

Evidence in repo: CQ15 returns no result, and the current final graph
does not yet expose a benchmark journey individual for the
Brixton-to-Canary-Wharf example.

Why it matters: journey modelling is present in the ontology, but the
instance layer does not yet carry the same benchmark-style coverage.

Completion strategy: derive representative journeys from a structured
routing source if available, or preserve a curated benchmark journey
layer.

\subsubsection{4. Planned engineering-closure
instances}\label{4-planned-engineering-closure-instances}

Evidence in repo: CQ13 returns no result, and the current artifact does
not expose \texttt{tfl:EngineeringClosure} instances in the final graph.

Why it matters: service-status information is part of the structured
pipeline, but a clearer class-materialization step would improve
disruption-oriented queries.

Completion strategy: materialize disruption individuals from validated
line-status payloads and connect them through
\texttt{tfl:hasDisruption}.

\subsubsection{5. Narrow operator answers for benchmark
queries}\label{5-narrow-operator-answers-for-benchmark-queries}

Evidence in repo: queries such as CQ8 and CQ14 return broad sets of
operator-like entities rather than a single narrow benchmark answer.

Why it matters: operator facts are present, but they would benefit from
a stricter mapping between organization names, line entities, and the
canonical operator class.

Completion strategy: introduce an ontology-aware operator-resolution
stage before final serialization.

\subsection{Retrieval-Augmented Completion
Strategy}\label{retrieval-augmented-completion-strategy}

The best completion strategy for the current graph is a KG-centred RAG
workflow that uses the existing graph as the first layer of control
rather than treating the model as a free-form generator.

\begin{enumerate}
\def\labelenumi{\arabic{enumi}.}
\tightlist
\item
  KG retrieval: retrieve the exact entities, classes, and predicates
  involved in missing or weak competency-question answers.
\item
  Source retrieval: retrieve only the supporting TfL cache records and
  relevant text passages for those entities.
\item
  Controlled completion: ask the model to propose additions only within
  the ontology\textquotesingle s class and property vocabulary.
\item
  Validation: accept completions only when they satisfy domain, range,
  entity-resolution, and benchmark-consistency checks.
\end{enumerate}

This approach is preferable because the graph already has substantial
recall. The main need is to improve precision and targeted completion
without weakening the automated pipeline.

\subsection{RAG Results On Three Worked
Gaps}\label{rag-results-on-three-worked-gaps}

\subsubsection{Case 1: Ontology-alignment
repair}\label{case-1-ontology-alignment-repair}

Retrieved gap: auxiliary runtime types are present alongside the
canonical coursework ontology classes.

RAG completion action:

\begin{itemize}
\tightlist
\item
  retrieve the auxiliary class labels and their connected entities
\item
  retrieve the corresponding canonical ontology classes from the TBox
\item
  ask the model to propose one-to-one alignments or justified
  declarations
\item
  validate the alignments before applying them
\end{itemize}

Expected result: a graph that is easier to query consistently and more
closely aligned with its own ontology.

\subsubsection{Case 2: Benchmark fare and journey
enrichment}\label{case-2-benchmark-fare-and-journey-enrichment}

Retrieved gap: the current graph lacks benchmark fare and journey
individuals required by selected coursework queries.

RAG completion action:

\begin{itemize}
\tightlist
\item
  retrieve the relevant fare, station, and route entities from the
  current graph
\item
  retrieve supporting structured records or narrow textual evidence
\item
  ask the model to propose only ontology-valid fare and journey
  instances
\item
  validate against the intended class and property structure before
  insertion
\end{itemize}

Expected result: stronger answerability for fare- and journey-oriented
competency questions without abandoning the automated architecture.

\subsubsection{Case 3: Precision-constrained night-service
refinement}\label{case-3-precision-constrained-night-service-refinement}

Retrieved gap: night-service facts are populated, but the benchmark
query for Night Tube lines needs a tighter answer set.

RAG completion action:

\begin{itemize}
\tightlist
\item
  retrieve all entities carrying night-service predicates
\item
  retrieve their current type assertions and line metadata
\item
  keep only entities validated as Underground lines
\item
  regenerate or filter the final Night Tube answer set from that
  controlled subset
\end{itemize}

Expected result: cleaner separation between Night Tube lines and night
bus services.

\subsection{Recommended Next Actions}\label{recommended-next-actions}

\begin{enumerate}
\def\labelenumi{\arabic{enumi}.}
\tightlist
\item
  Add a post-build ontology-alignment pass for auxiliary runtime
  classes.
\item
  Materialize benchmark fare, journey, disruption, and facility
  instances from validated evidence.
\item
  Tighten operator and night-service entity resolution before final
  serialization.
\item
  Keep the source inventory and prompt documentation synchronized with
  each pipeline revision.
\item
  Use a KG-centred RAG workflow to improve targeted completion without
  reducing reproducibility.
\end{enumerate}

\subsection{Conclusion}\label{conclusion}

The current graph is already substantial in scale, provenance, and
automation. Its completion work is therefore best understood as targeted
enrichment: aligning auxiliary classes with the ontology, materializing
a small number of high-value benchmark instances, and refining selected
query-sensitive entity types. This keeps the project faithful to the
coursework brief while providing a clear path to stronger
competency-question coverage.

\end{document}
